% Generated draft: requires mathematical and source review.
\documentclass[11pt]{article}
\usepackage[margin=25mm]{geometry}
\usepackage{fontspec}
\setmainfont{Latin Modern Roman}
\usepackage{amsmath,amssymb,mathtools,unicode-math}
\setmathfont{Latin Modern Math}
\usepackage{enumitem,needspace,xurl,hyperref}
\setlength{\parindent}{0pt}
\setlength{\parskip}{5pt}
\setlength{\emergencystretch}{3em}
\newcommand{\N}{\mathbb{N}}
\newcommand{\Z}{\mathbb{Z}}
\newcommand{\Q}{\mathbb{Q}}
\newcommand{\R}{\mathbb{R}}
\newcommand{\C}{\mathbb{C}}
\newcommand{\F}{\mathbb{F}}
\newcommand{\A}{\mathbb{A}}
\newcommand{\PP}{\mathbb{P}}
\newcommand{\E}{\mathbb{E}}
\newcommand{\eps}{\varepsilon}
\newcommand{\dd}{\,\mathrm{d}}
\begin{document}
\section*{UnsolvedMath ID: 30001605}
\textbf{Source problem code:} catalogue entry
\subsection{Definitions and mathematical statement}
Let \(f\colon X\to B\) be a morphism of complex algebraic varieties. A morphism is proper if it is separated, of finite type, and universally closed. For \(b\in B(\C)\), write \(X_b=X\times_B\operatorname{Spec}\C\) for its fiber. A smooth, projective, connected complex curve \(C\) has genus one if \(\dim_{\C}H^1(C,\mathcal O_C)=1\). In pointed terminology, an elliptic curve additionally has a distinguished point \(O\in C(\C)\).

The supplied record describes the workshop's common convention as follows: an elliptic fibration is a proper morphism of complex algebraic varieties whose general fiber is an elliptic curve. If ``general fiber'' means fibers over a nonempty dense Zariski-open subset, and ``elliptic curve'' here refers to an underlying genus-one curve without a specified origin, this convention has the formulation
\[
 f\colon X\to B\text{ is proper},\qquad
 \exists\,U\subseteq B\text{ nonempty dense Zariski-open}\quad
 \forall b\in U(\C),\quad
 X_b\text{ is a smooth, projective, connected curve of genus }1.
\]
This is an explicit interpretation of a terminology fragment, not a recovered conjecture. The record does not specify its conventions for varieties, the precise meaning of general fiber, or whether fiberwise origins are part of the data. In particular, it does not require a section \(s\colon B\to X\) satisfying \(f\circ s=\operatorname{id}_B\). Smoothness, existence of a section, flatness, and prescribed singular fibers are listed as possible refinements, without selecting any of them; the intended smoothness requirement and any restrictions on fibers outside \(U\) remain unspecified. No determinate open mathematical assertion is supplied.
\subsection{Short English statement}
The record discusses the meaning of an elliptic fibration: a proper morphism of complex algebraic varieties with general fiber an elliptic curve. Additional requirements vary with the convention. The supplied text is a definition discussion, so it does not determine an open problem to formalize.
\subsection{Sources}
\begin{itemize}
\item Local UnsolvedMath registry, supplied record with permanent ID \(30001605\), title ``Defining Higher-Dimensional Elliptic Fibrations,'' exact source code \texttt{catalogue entry}, in the Oberwolfach Reports Open Problems collection. No individual proposer is supplied; the proposed year is recorded as 2010.
\item Supplied source citation: \emph{Mini-Workshop: Higher Dimensional Elliptic Fibrations} (2011), Oberwolfach Reports, pp.\ 2651--2680, DOI \texttt{10.4171/owr/2010/46}. Source URL and supplied literature-source link: \url{https://doi.org/10.4171/owr/2010/46}.
\item Literature assessment dated exactly \texttt{2026-08-21}: supplied dated curation, not independently verified. It retains the label ``open'' because the extracted sentence is a definitional fragment rather than a precise conjecture and therefore cannot be verified as solved under the stated curation rule. The supplied background requests independent verification of time-sensitive status and cited claims.
\item The supplied statement-verification field describes the report section ``What is an elliptic fibration?'' as a terminology discussion with optional refinements, not a determinate open problem. The registry marks the statement status as \texttt{unrecoverable}, supplies no clean statement, and retains the original record for provenance. These are supplied source qualifications, not claims of independent source inspection.
\end{itemize}
\end{document}
